%=============================================================================
% WAVE 3, ITERATION 5: PATENT 1 UPDATE
% Field Drag / Gravitational Cloaking ---
% M_psi Correction Propagation, T1 Falsification Impact,
% T5 Source-Term Inconsistency, Corrected Smoking-Gun Specification
%
% Agent 1 | DFD Research Programme
% Based on: W3i4_P1_update.tex, W3i4_P11_update.tex (M_psi correction),
%           W3i4_P14_update.tex (T1 falsification, T5 dual-role),
%           ITERATION_TRACKER.md (W3i4 summary)
% Date: 20 March 2026
%=============================================================================

%--- Preamble (section-only, no \documentclass per specification) ----------
\usepackage{amsmath,amssymb,amsfonts,amsthm,mathtools}
\usepackage{booktabs}
\usepackage{longtable}
\usepackage{array}
\usepackage{colortbl}
\usepackage{xcolor}
\usepackage{enumitem}
\usepackage{siunitx}
\usepackage{hyperref}
\usepackage{tcolorbox}
\tcbuselibrary{skins,breakable}

%--- Macros -----------------------------------------------------------------
\newcommand{\dpsi}{\delta\psi}
\newcommand{\lam}{\lambda}
\newcommand{\SNR}{\mathrm{SNR}}
\newcommand{\order}[1]{\mathcal{O}(#1)}
\newcommand{\half}{\tfrac{1}{2}}
\newcommand{\ee}[1]{\times 10^{#1}}
\newcommand{\ns}{\,\mathrm{ns}}
\newcommand{\ps}{\,\mathrm{ps}}
\newcommand{\fs}{\,\mathrm{fs}}
\newcommand{\muHz}{\,\mu\mathrm{Hz}}
\newcommand{\lbone}{|\lam - 1|}
\newcommand{\Meff}{M_\psi^{\rm eff}}
\newcommand{\Mpsiused}{M_\psi^{\rm used}}
\newcommand{\psicosmo}{\bar{\psi}_{\rm cosmo}}

\definecolor{strongcolor}{rgb}{0,0.45,0}
\definecolor{weakcolor}{rgb}{0.65,0,0}
\definecolor{conjcolor}{rgb}{0.6,0.3,0}
\definecolor{rowhi}{rgb}{0.85,0.95,0.85}
\definecolor{rowmed}{rgb}{0.95,0.95,0.80}
\definecolor{rowlo}{rgb}{0.97,0.87,0.87}
\definecolor{falsified}{rgb}{0.7,0,0}

\newtheorem{theorem}{Theorem}[section]
\newtheorem{proposition}[theorem]{Proposition}
\newtheorem{corollary}[theorem]{Corollary}
\newtheorem{lemma}[theorem]{Lemma}
\newtheorem{definition}[theorem]{Definition}
\newtheorem{finding}[theorem]{Finding}
\newtheorem{correction}[theorem]{Correction}

%=============================================================================
\section{W3i5 P1: M$_\psi$ Correction Propagation, T1/T5 Impact, Corrected
Smoking-Gun}
\label{sec:w3i5_intro}
%=============================================================================

\noindent\textbf{Mandate.}
W3i4 established three critical results that jointly threaten the entire
P1 experimental programme:
\begin{enumerate}[nosep]
\item \textbf{M$_\psi$ correction} (W3i4-P11, Correction~3.1):
  $\Mpsiused = 10^{-6}$~kg was a placeholder; the correct effective inertial
  mass is $\Meff = 3.01\ee{6}$~kg (TESLA cavity) or
  $7.84\ee{5}$~kg (sub-litre re-entrant). All \emph{absolute} SNR estimates
  that depended on $M_\psi$ in the denominator were $10^6\times$
  over-optimistic.
\item \textbf{T1 falsification} (W3i4-P14, Section~6): The minimal DFD
  cosmological coupling $G_{\rm eff}(t) = G_N e^{-2\psicosmo(t)}$ predicts
  $\dot{G}/G = +2.6\ee{-11}$~yr$^{-1}$, exceeding the LLR bound by
  $347\times$. Declared genuine falsification of the minimal cosmological
  sector.
\item \textbf{T5 source-term inconsistency} (W3i4-P14, Section~4): The DFD
  psi-field equation is sourced only by EM stress-energy
  ($\Box\psi \propto T_{\rm EM}$), yet the Mach integral treats $\psi$ as the
  Newtonian gravitational potential ($\psi = GM/c^2 r$, sourced by matter
  density). These two roles are mutually inconsistent.
\end{enumerate}

This W3i5 update addresses four questions:
\begin{enumerate}[nosep,label=\textbf{Q\arabic*.}]
\item Does the $10^{6}\times$ M$_\psi$ correction affect the 4.0~ns bubble
  timing signal?
\item What survives after honest propagation of all corrections?
\item How do T1 and T5 affect P1's theoretical foundation?
\item What is the corrected smoking-gun specification?
\end{enumerate}

%=============================================================================
\section{Q1: Does the M$_\psi$ Correction Affect the 4.0 ns Signal?}
\label{sec:mpsi_propagation}
%=============================================================================

\subsection{Tracing M$_\psi$ through the P1 signal chain}
\label{ssec:mpsi_trace}

The W3i4 P1 signal model (Eq.~(1)--(4) of W3i4-P1) gives:
\begin{equation}
\delta t_{\rm NR} = \frac{2L\,|\nabla\dpsi|\sin\theta}{c}
= \frac{2L\,|\dpsi_{\rm bubble}|}{c\,R_b}.
\label{eq:signal_recap}
\end{equation}

The critical question is: \emph{where does $M_\psi$ enter this expression?}

The signal depends on three quantities:
\begin{enumerate}[nosep]
\item $L$ (baseline length) --- geometric, no $M_\psi$ dependence.
\item $R_b$ (bubble radius) --- geometric, set by cavity array geometry.
\item $|\dpsi_{\rm bubble}|$ --- the psi-field amplitude inside the bubble.
\end{enumerate}

The psi-field amplitude is determined by the DFD source equation. In the
static, linearised regime:
\begin{equation}
\nabla^2 \dpsi = -\frac{(\lam - 1)}{c^2}\,u_{\rm EM}(\mathbf{r}),
\label{eq:source_eq}
\end{equation}
where $u_{\rm EM}$ is the EM energy density (J/m$^3$). The solution:
\begin{equation}
|\dpsi_{\rm bubble}| \sim \frac{(\lam - 1)\,u_{\rm EM}\,R_b^2}{c^2}.
\label{eq:dpsi_from_source}
\end{equation}

\begin{theorem}[$M_\psi$ Does Not Appear in the P1 Signal Formula]
\label{thm:mpsi_absent}
The nonreciprocal timing signal $\delta t_{\rm NR}$ depends on
$|\dpsi_{\rm bubble}|$, which is determined by the static Poisson equation
(\ref{eq:source_eq}). The effective mode mass $\Meff$ enters only in the
\emph{dynamical} response of the psi-mode --- specifically, the parametric
oscillation amplitude and the SQUID readout signal in P11/P2 cavity
experiments. It does \emph{not} enter the static or quasi-static psi-field
amplitude that determines the bubble strength.

Specifically:
\begin{itemize}[nosep]
\item The static solution $\dpsi \propto (\lam-1)\,u_{\rm EM}/c^2$ contains
  no mass parameter.
\item The mode mass $\Meff = \mu_0^{\rm gal}\,V_{\rm cav}/(4\pi G)$ governs
  the oscillation frequency $\Omega_\psi$ and the SQUID-coupled displacement
  sensitivity $\delta q \propto 1/\sqrt{\Meff}$, which are relevant only
  for \emph{resonant detection} experiments (P11 cavity readout, P2
  parametric amplification of the \emph{mode amplitude}).
\item The P1+P5 timing experiment measures the \emph{accumulated phase delay}
  through a macroscopic psi-bubble, not the parametric response of a single
  cavity mode. The measurement chain is: EM source $\to$ static $\dpsi$ field
  $\to$ refractive-index gradient $\to$ one-way light travel time.
\end{itemize}
\end{theorem}

\begin{tcolorbox}[colback=green!5!white, colframe=green!50!black,
  title=\textbf{Result: The 4.0 ns Signal Is Unaffected by the M$_\psi$ Correction}]
The W3i4 signal value $\delta t_{\rm NR} = 4.0$~ns stands. The $10^6\times$
over-optimism identified in W3i4-P11 applies to:
\begin{itemize}[nosep]
\item SQUID-active absolute sensitivity projections (P11 single-cavity,
  P2+P11 hybrid, raw SQUID PSD calculations).
\item Any calculation where $M_\psi$ appears in the denominator of a
  signal amplitude.
\end{itemize}
It does \textbf{not} apply to:
\begin{itemize}[nosep]
\item The static psi-bubble amplitude (P1 operating point).
\item The P1+P5 nonreciprocal timing signal.
\item The SNR baseline of 4000 (which is $\delta t_{\rm NR}/\sigma_{\rm clock}$,
  neither of which involves $M_\psi$).
\item Any of the W3i4 P1 configuration SNR estimates (Alpha, Beta, Gamma).
\end{itemize}
\end{tcolorbox}

\subsection{Verification: Dimensional and physical consistency check}
\label{ssec:verification}

To be rigorous, we verify that $M_\psi$ cancels explicitly in the P1 signal
chain.

The psi-bubble at the P1 operating point requires $|\dpsi| = 0.30$. From
Eq.~(\ref{eq:dpsi_from_source}):
\begin{equation}
|\dpsi| = 0.30 \implies u_{\rm EM} = \frac{0.30\,c^2}{(\lam-1)\,R_b^2}.
\label{eq:uEM_required}
\end{equation}

For $\lbone = 1.85\ee{-2}$ and $R_b = 5$~m:
\begin{equation}
u_{\rm EM} = \frac{0.30 \times 9\ee{16}}{1.85\ee{-2} \times 25}
= \frac{2.7\ee{16}}{0.4625} = 5.84\ee{16}~\text{J/m}^3.
\label{eq:uEM_numerical}
\end{equation}

This is an enormous energy density ($\sim 10^{16}$~J/m$^3$, comparable to
nuclear energy densities). The SRF cavity with $Q = 10^{10}$ and feed power
$P_{\rm feed}$ stores energy $E_{\rm stored} = Q\,P_{\rm feed}/\omega_{\rm SRF}$.
For 35~GHz ($\omega = 2.2\ee{11}$~rad/s) and $P_{\rm feed} = 1$~kW:
\begin{equation}
E_{\rm stored} = \frac{10^{10} \times 10^3}{2.2\ee{11}} = 45.5~\text{J}.
\label{eq:E_stored}
\end{equation}

In a TESLA cavity volume $V_{\rm cav} = 3.84\ee{-3}$~m$^3$:
\begin{equation}
u_{\rm EM}^{(\rm SRF)} = \frac{45.5}{3.84\ee{-3}} = 1.18\ee{4}~\text{J/m}^3.
\label{eq:uEM_SRF}
\end{equation}

\begin{correction}[Required Energy Density vs.\ Available: 12-Order Gap]
\label{corr:energy_gap}
The psi-bubble operating point $|\dpsi| = 0.30$ requires
$u_{\rm EM} = 5.84\ee{16}$~J/m$^3$ (Eq.~\ref{eq:uEM_numerical}), while the
best achievable SRF cavity provides
$u_{\rm EM}^{(\rm SRF)} = 1.18\ee{4}$~J/m$^3$ (Eq.~\ref{eq:uEM_SRF}).

The ratio:
\begin{equation}
\frac{u_{\rm EM}^{(\rm required)}}{u_{\rm EM}^{(\rm SRF)}}
= \frac{5.84\ee{16}}{1.18\ee{4}} = 4.95\ee{12}.
\label{eq:energy_gap}
\end{equation}

\textbf{The required EM energy density exceeds the available by a factor of
$\sim 5\ee{12}$.}

This means the actual psi-bubble amplitude achievable with a TESLA-class SRF
cavity at 1~kW feed power is:
\begin{equation}
|\dpsi|_{\rm actual} = 0.30 \times \frac{u_{\rm EM}^{(\rm SRF)}}{u_{\rm EM}^{(\rm required)}}
= 0.30 \times \frac{1.18\ee{4}}{5.84\ee{16}}
= 6.06\ee{-14}.
\label{eq:dpsi_actual}
\end{equation}
\end{correction}

\begin{tcolorbox}[colback=red!5!white, colframe=red!70!black,
  title=\textbf{CRITICAL FINDING: The 4.0 ns Signal Collapses --- Not From M$_\psi$
  But From Energy Density}]
The M$_\psi$ correction does not directly enter the P1 signal formula.
However, \emph{re-examining the signal from first principles} reveals a far
more devastating problem: the W3i4 analysis assumed the psi-bubble was at its
\emph{drag-reduction operating point} $|\dpsi| = 0.30$ without checking
whether this amplitude is physically achievable.

At the drag-reduction coupling $\lbone = 1.85\ee{-2}$ and 1~kW SRF feed power:
\begin{equation}
|\dpsi|_{\rm actual} = \frac{(\lam-1)\,u_{\rm EM}^{(\rm SRF)}\,R_b^2}{c^2}
= \frac{1.85\ee{-2} \times 1.18\ee{4} \times 25}{9\ee{16}}
= 6.06\ee{-14}.
\label{eq:dpsi_actual_box}
\end{equation}

The actual nonreciprocal timing signal:
\begin{equation}
\boxed{
\delta t_{\rm NR}^{(\rm actual)} = 4.0~\text{ns} \times \frac{6.06\ee{-14}}{0.30}
= 4.0~\text{ns} \times 2.02\ee{-13} = 8.1\ee{-22}~\text{s}
= 0.81~\text{zs (zeptoseconds)}.
}
\label{eq:signal_actual}
\end{equation}

The actual single-shot SNR:
\begin{equation}
\boxed{
\SNR_{\rm actual} = \frac{8.1\ee{-22}~\text{s}}{1.02\ee{-12}~\text{s}}
= 7.9\ee{-10} \approx 10^{-9}.
}
\label{eq:SNR_actual}
\end{equation}

The signal is nine orders of magnitude below the noise floor.
\end{tcolorbox}

\subsection{How did W3i3/W3i4 arrive at $|\dpsi| = 0.30$?}
\label{ssec:how_030}

The W3i1 analysis defined the ``drag-reduction operating point'' as the
psi-field amplitude required for 1\% drag reduction in a bluff body:
$|\dpsi|_{\rm drag} > 3.35\ee{-3}$, which maps to
$\lbone > 1.85\ee{-2}$. The operating point $|\dpsi| = 0.30$ was then stated
as a target amplitude, and subsequent iterations \emph{assumed} this value
was achieved without tracing back to the EM energy density required to
generate it.

The implicit assumption was that the SRF cavity could produce an arbitrarily
large $u_{\rm EM}$ by increasing $Q$ and/or $P_{\rm feed}$. However:

\begin{finding}[The $|\dpsi| = 0.30$ Operating Point Requires Nuclear-Scale
Energy Densities]
\label{find:nuclear}
At $\lbone = 1.85\ee{-2}$, achieving $|\dpsi| = 0.30$ in a 5~m radius
bubble requires $u_{\rm EM} = 5.84\ee{16}$~J/m$^3$. For comparison:
\begin{itemize}[nosep]
\item TNT detonation: $u \sim 10^{9}$~J/m$^3$
\item Nuclear reactor core: $u \sim 10^{10}$~J/m$^3$
\item Interior of a neutron star: $u \sim 10^{34}$~J/m$^3$
\item Highest laboratory EM field (MAGNET, pulsed): $B \sim 1200$~T,
  $u = B^2/(2\mu_0) \sim 5.7\ee{8}$~J/m$^3$
\item NIF laser focus (transient): $u \sim 10^{15}$~J/m$^3$
\end{itemize}

The required energy density is $\sim 60\times$ the NIF laser focus ---
the highest EM energy density ever achieved in a laboratory, and that
only for picosecond durations in a micron-scale focal volume. Sustaining
$5.8\ee{16}$~J/m$^3$ in a macroscopic volume is physically impossible with
any known or foreseeable technology.
\end{finding}

\begin{correction}[W3i1--W3i4 Assumed Unphysical Operating Point]
\label{corr:unphysical_operating_point}
The entire W3i1--W3i4 P1 analysis chain assumed the bubble operates at
$|\dpsi| = 0.30$ without verifying that this amplitude can be generated by
any physical EM source. This assumption is wrong by $\sim 12$ orders of
magnitude. The correct procedure is to compute $|\dpsi|$ from the available
$u_{\rm EM}$ and $\lbone$, then derive the timing signal from that.

\textbf{All W3i1--W3i4 results that assumed $|\dpsi| = 0.30$ are invalid
as physical predictions, though they remain valid as conditional statements:
``\emph{if} $|\dpsi| = 0.30$ could be achieved, \emph{then} the signal would
be 4.0~ns.''}
\end{correction}

%=============================================================================
\section{Q2: What Survives? --- Honest Signal Estimates from First Principles}
\label{sec:honest_signal}
%=============================================================================

\subsection{Signal formula from achievable parameters}
\label{ssec:honest_formula}

Starting from the DFD source equation and working forward:
\begin{equation}
\delta t_{\rm NR} = \frac{2L}{c} \cdot \frac{|\dpsi|_{\rm actual}}{R_b}
= \frac{2L}{c} \cdot \frac{(\lam-1)\,u_{\rm EM}\,R_b}{c^2}.
\label{eq:honest_signal}
\end{equation}

\textbf{Parameters under experimental control:}
\begin{itemize}[nosep]
\item $L$: baseline length (can be 1--100~m in a lab)
\item $u_{\rm EM}$: EM energy density in the source region
\item $R_b$: bubble radius (set by source geometry; for a cavity,
  $R_b \sim V_{\rm cav}^{1/3}$)
\item $\lbone$: \emph{not} under control; this is the unknown coupling constant.
\end{itemize}

\textbf{Best achievable EM parameters:}
\begin{itemize}[nosep]
\item TESLA SRF cavity: $Q = 2\ee{10}$, $V_{\rm cav} = 3.84\ee{-3}$~m$^3$,
  $R_b \approx 0.16$~m (cavity half-length).
\item Feed power (CW, sustainable): $P_{\rm feed} = 10$~kW.
\item $E_{\rm stored} = Q\,P_{\rm feed}/\omega = 2\ee{10} \times 10^4/(2.2\ee{11})
  = 909$~J.
\item $u_{\rm EM} = 909/3.84\ee{-3} = 2.37\ee{5}$~J/m$^3$.
\end{itemize}

\textbf{Alternative: pulsed high-field magnet:}
\begin{itemize}[nosep]
\item 100~T field, $V = 10^{-3}$~m$^3$ ($\ell = 0.1$~m bore, 0.1~m long):
  $u_{\rm EM} = B^2/(2\mu_0) = (100)^2/(2 \times 1.26\ee{-6}) = 3.98\ee{9}$~J/m$^3$.
\item $R_b \approx 0.05$~m.
\end{itemize}

\textbf{Alternative: NIF-class pulsed laser (transient):}
\begin{itemize}[nosep]
\item $P_{\rm peak} = 500$~TW, focal volume $V \sim 10^{-12}$~m$^3$:
  $u_{\rm EM} \sim 5\ee{14}$~J/m$^3$.
\item $R_b \sim 10^{-4}$~m. Duration $\sim 10^{-11}$~s.
\end{itemize}

\subsection{Timing signal for each source configuration}
\label{ssec:timing_configs}

Using Eq.~(\ref{eq:honest_signal}) with $L = R_b$ (optimal geometry:
baseline crosses the source region):

\begin{equation}
\delta t_{\rm NR} = \frac{2\,(\lam-1)\,u_{\rm EM}\,R_b^2}{c^3}.
\label{eq:honest_signal_simplified}
\end{equation}

\begin{center}
\textbf{Table 1: Honest P1 Timing Signal at Various $\lbone$ Values}
\vspace{0.5em}

\begin{tabular}{@{}llcccc@{}}
\toprule
\textbf{Source} & $u_{\rm EM}$ (J/m$^3$) & $R_b$ (m) &
$\lbone = 10^{-2}$ & $\lbone = 10^{-6}$ & $\lbone = 10^{-9}$ \\
\midrule
\multicolumn{3}{@{}l}{\emph{Signal $\delta t_{\rm NR}$ (seconds):}} \\
\midrule
SRF (10 kW) & $2.37\ee{5}$ & 0.16 & $2.24\ee{-20}$ & $2.24\ee{-24}$
  & $2.24\ee{-27}$ \\
Pulsed magnet (100 T) & $3.98\ee{9}$ & 0.05 & $7.37\ee{-20}$ & $7.37\ee{-24}$
  & $7.37\ee{-27}$ \\
NIF laser (transient) & $5\ee{14}$ & $10^{-4}$ & $3.70\ee{-23}$
  & $3.70\ee{-27}$ & $3.70\ee{-30}$ \\
\midrule
\multicolumn{3}{@{}l}{\emph{SNR (single-shot, $\sigma = 1$~ps):}} \\
\midrule
SRF (10 kW) & & & $2.2\ee{-8}$ & $2.2\ee{-12}$ & $2.2\ee{-15}$ \\
Pulsed magnet (100 T) & & & $7.4\ee{-8}$ & $7.4\ee{-12}$ & $7.4\ee{-15}$ \\
NIF laser (transient) & & & $3.7\ee{-11}$ & $3.7\ee{-15}$ & $3.7\ee{-18}$ \\
\bottomrule
\end{tabular}
\end{center}

\begin{finding}[No Laboratory Configuration Achieves SNR $\geq 1$]
\label{find:no_detection}
Even at $\lbone = 10^{-2}$ (four orders above the passive bound of
$4.9\ee{-7}$, and completely excluded by cavity Q data), the best
laboratory source (pulsed 100~T magnet) achieves
$\SNR = 7.4\ee{-8}$. The signal is 7 orders of magnitude below
the noise floor.

At the SQMS hint level $\lbone \sim 10^{-6}$:
$\SNR_{\rm best} = 7.4\ee{-12}$ (12 orders below detection).

At the current passive bound $\lbone < 1.56\ee{-9}$:
$\SNR_{\rm best} = 7.4\ee{-15}$ (15 orders below detection).

\textbf{The P1+P5 bubble timing experiment cannot detect the DFD psi-field
with any physically realisable EM source at any currently viable coupling
constant.}
\end{finding}

\subsection{What averaging could help?}
\label{ssec:averaging}

The SNR improves as $\sqrt{N_{\rm rep}}$. To reach SNR $= 5$ from the best
configuration (pulsed magnet at $\lbone = 10^{-2}$, $\SNR_1 = 7.4\ee{-8}$):
\begin{equation}
N_{\rm rep} = \left(\frac{5}{7.4\ee{-8}}\right)^2 = 4.6\ee{15}.
\label{eq:reps_needed}
\end{equation}

At 1 measurement per second: $4.6\ee{15}$~s $= 1.4\ee{8}$~years.

At the SQMS hint level ($\lbone = 10^{-6}$): $N_{\rm rep} = 4.6\ee{23}$,
requiring $1.4\ee{16}$~years (longer than the age of the universe).

\begin{finding}[Averaging Cannot Rescue the Experiment]
\label{find:averaging_fails}
No physically realisable measurement repetition rate can compensate for the
7--15 order-of-magnitude SNR deficit. The fundamental problem is the tiny
$|\dpsi|$ achievable with laboratory EM fields, not the noise floor
(which is already excellent at 1~ps).
\end{finding}

\subsection{Re-examining the W3i4 configurations}
\label{ssec:configs_corrected}

\begin{center}
\textbf{Table 2: W3i4 Configurations --- Corrected with Actual $|\dpsi|$}
\vspace{0.5em}

\begin{tabular}{@{}lcccc@{}}
\toprule
& \textbf{W3i4 (assumed)} & \textbf{W3i5 (corrected)} & \textbf{Ratio} \\
\midrule
$|\dpsi|_{\rm bubble}$ & 0.30 & $6.1\ee{-14}$ & $4.9\ee{12}$ \\
$\delta t_{\rm NR}$ & 4.0~ns & 0.81~zs & $4.9\ee{12}$ \\
$\SNR_{\rm Baseline}$ & 4000 & $7.9\ee{-10}$ & $5.0\ee{12}$ \\
$\SNR_{\rm Alpha}$ & $1.3\ee{5}$ & $2.6\ee{-8}$ & $5.0\ee{12}$ \\
$\SNR_{\rm Beta}$ & $1.1\ee{5}$ & $2.2\ee{-8}$ & $5.0\ee{12}$ \\
$\SNR_{\rm Gamma}$ & $3.4\ee{6}$ & $6.8\ee{-7}$ & $5.0\ee{12}$ \\
\bottomrule
\end{tabular}
\end{center}

\textbf{All configurations are undetectable.} The multi-patent gain factors
(3.3$\times$, 10$\times$, etc.) are real but irrelevant: they multiply a
signal that is already $10^{12}$ below the noise.

\subsection{Could a different geometry help?}
\label{ssec:different_geometry}

The signal scales as $u_{\rm EM} \cdot R_b^2$. This means the optimal
strategy is to maximise the product $u_{\rm EM} \cdot R_b^2$, which is
proportional to the \emph{total EM energy} in the source:
$E_{\rm total} = u_{\rm EM} \cdot V \sim u_{\rm EM} \cdot R_b^3$. Thus:
\begin{equation}
\delta t_{\rm NR} \propto E_{\rm total} / R_b
\propto E_{\rm total}^{2/3} \cdot u_{\rm EM}^{1/3}.
\label{eq:scaling_geometry}
\end{equation}

This favours high total stored energy over high energy density. The
world's largest superconducting magnet energy storage:
\begin{itemize}[nosep]
\item LHC dipole circuit: $E_{\rm stored} = 10$~GJ at $B = 8.3$~T,
  $V \sim 10^4$~m$^3$.
\item $u_{\rm EM} = B^2/(2\mu_0) = (8.3)^2/(2.52\ee{-6}) = 2.74\ee{7}$~J/m$^3$.
\item $R_b \approx V^{1/3} \approx 21.5$~m.
\end{itemize}

The signal from the LHC dipole circuit (if the entire stored field were
used as a psi-source):
\begin{equation}
\delta t_{\rm NR}^{(\rm LHC)} = \frac{2 \times (\lam-1) \times 2.74\ee{7}
\times (21.5)^2}{(3\ee{8})^3}
= \frac{2 \times (\lam-1) \times 1.27\ee{10}}{2.7\ee{25}}.
\label{eq:LHC_signal}
\end{equation}

At $\lbone = 10^{-2}$:
$\delta t_{\rm NR}^{(\rm LHC)} = 2 \times 10^{-2} \times 1.27\ee{10}/(2.7\ee{25})
= 9.4\ee{-18}$~s $= 9.4$~as.

$\SNR_{\rm LHC}^{(\lbone=10^{-2})} = 9.4\ee{-18}/(1.0\ee{-12}) = 9.4\ee{-6}$.

Still six orders below detection. At $\lbone = 10^{-6}$: $\SNR = 9.4\ee{-10}$.

\begin{theorem}[No Known or Foreseeable EM Source Can Detect the P1 Signal]
\label{thm:no_source}
Even using the entire LHC stored energy (10~GJ, the largest concentrated
EM energy on Earth) as a DFD psi-source, and even assuming a coupling
constant four orders above the passive bound ($\lbone = 10^{-2}$), the
nonreciprocal timing signal is $\SNR \approx 10^{-5}$.

At the current passive bound $\lbone < 1.56\ee{-9}$:
$\SNR \approx 10^{-12}$.

\textbf{The P1+P5 nonreciprocal timing experiment is unfeasible.}
\end{theorem}

%=============================================================================
\section{What Went Wrong in W3i1--W3i4?}
\label{sec:postmortem}
%=============================================================================

\subsection{Root cause analysis}
\label{ssec:root_cause}

The error chain:
\begin{enumerate}[nosep]
\item W3i1 defined $|\dpsi|_{\rm drag} > 3.35\ee{-3}$ as the threshold for
  measurable drag and $\lbone > 1.85\ee{-2}$ as the required coupling. This
  was correct.
\item W3i1 then set the \emph{operating point} at $|\dpsi| = 0.30$
  (100$\times$ the threshold), without computing the EM energy density
  required to generate this field.
\item W3i2 and W3i3 inherited $|\dpsi| = 0.30$ as a given parameter.
\item W3i4 computed the full noise budget (correct), the 15-patent
  interaction matrix (correct conditional on the signal), and the
  smoking-gun specification (correct conditional on the signal), but
  never re-derived $|\dpsi|$ from first principles.
\item The M$_\psi$ correction (W3i4-P11) was correctly identified as
  $10^6\times$ over-optimistic for SQUID readout, but this was a
  \emph{different} problem from the P1 signal issue.
\end{enumerate}

\textbf{The root cause is that $|\dpsi| = 0.30$ was treated as an
input parameter rather than a derived quantity.} When derived from
achievable EM energy densities and the DFD coupling, the actual bubble
amplitude is $\sim 10^{-13}$, not $0.30$.

\subsection{Are the relative results from W3i1--W3i4 still valid?}
\label{ssec:relative_results}

Yes, the following survive as correct \emph{relative} statements:
\begin{itemize}[nosep]
\item The drag threshold requires $\lbone > 1.85\ee{-2}$ (W3i1).
\item K-H stabilisation threshold is five orders lower (W3i1).
\item DFD PPN parameters $\beta = \gamma = 1$ (W3i4).
\item The 15-patent interaction matrix (W3i4): relative gains $G_k$ are
  correct.
\item The noise budget: 96.1\% clock, 3.9\% EM source (W3i4).
\item The NS timing anomaly: 2.15\%, 6.0~ps (W3i4 correction).
\item The non-reciprocal drag geometry (W3i4).
\end{itemize}

The absolute predictions ($\delta t = 4.0$~ns, $\SNR = 4000$, smoking-gun
specification) are \textbf{invalidated as achievable experimental targets}.

%=============================================================================
\section{Q3: Impact of T1 Falsification and T5 Inconsistency on P1}
\label{sec:T1_T5_impact}
%=============================================================================

\subsection{T1 impact on P1}
\label{ssec:T1_impact}

T1 states: the minimal DFD cosmological coupling predicts
$\dot{G}/G = +2.6\ee{-11}$~yr$^{-1}$, exceeding the LLR bound by
$347\times$.

\begin{finding}[T1 Does Not Directly Affect P1 Laboratory Predictions]
\label{find:T1_P1}
The P1 experiment operates in the \emph{laboratory EM-coupling sector}
of DFD, not the cosmological sector. The relevant P1 equations are:
\begin{enumerate}[nosep]
\item The psi-field source equation: $\nabla^2\dpsi = -(\lam-1)\,u_{\rm EM}/c^2$.
\item The refractive-index formula: $n = e^{\dpsi}$.
\item The timing delay: $\delta t = 2L|\nabla\dpsi|/c$.
\end{enumerate}

None of these involve $\psicosmo$, $\dot{G}/G$, or the Mach integral.
The T1 falsification concerns the \emph{cosmological background} sector
($\psicosmo \ne 0$ with evolving $G_{\rm eff}$), which is decoupled from
the laboratory EM-sourced $\dpsi$ at the level of the current DFD
field equations.
\end{finding}

\textbf{However}, T1 has an important \emph{indirect} impact:

\begin{finding}[T1 Undermines the Theoretical Motivation for P1]
\label{find:T1_indirect}
The DFD theory claims that $\psi$ is a universal scalar field that:
(a) mediates gravity at all scales,
(b) is sourced by EM stress-energy at the laboratory scale, and
(c) has a cosmological background $\psicosmo$ that produces MOND-like
dynamics and a time-varying $G$.

T1 falsifies claim~(c). If the cosmological sector is wrong, the
theoretical unification that connects laboratory EM-coupling to
astrophysical observations is broken. The laboratory predictions
(drag reduction, timing signal) remain mathematically self-consistent
but lose their connection to astrophysical evidence.

Specifically: the MOND-fit that constrains $\mu_0 = 0.657$ relies on the
same $\psi$-as-gravitational-potential identification that T1 falsifies.
If $\psi$ is \emph{not} the gravitational potential (as T5 suggests), then
$\mu_0$ is not constrained by galactic rotation curves, and the coupling
parameter $\lbone$ has no astrophysical anchor.
\end{finding}

\subsection{T5 impact on P1}
\label{ssec:T5_impact}

T5 (from W3i4-P14, Section~4) states: the DFD psi-field has two
mutually inconsistent roles:
\begin{enumerate}[nosep]
\item \textbf{Gravitational role}: $\psi = GM/(c^2 r)$, sourced by matter
  density $\rho_{\rm matter}$.
\item \textbf{Field-equation role}: $\Box\psi = (4\pi G/c^4)\,T_{\rm EM}$,
  sourced only by EM stress-energy.
\end{enumerate}

\begin{theorem}[T5 Creates a Dilemma for P1]
\label{thm:T5_dilemma}
P1 requires the psi-field to be sourced by EM energy (role~2) to generate
a laboratory psi-bubble. This is consistent with the field equation.

However, the physical mechanism by which $\dpsi$ creates a refractive
index (and hence a timing delay) assumes $\psi$ acts as a gravitational
potential (role~1): the metric $g_{\mu\nu} = \text{diag}(-e^{-\dpsi},
e^{+\dpsi}, \ldots)$ is motivated by the identification of $\psi$ with
the Newtonian potential.

\textbf{If role~1 and role~2 are inconsistent (T5), then the same $\dpsi$
generated by an EM source cannot be assumed to curve spacetime as a
gravitational potential would.}

Three possibilities:
\begin{enumerate}[nosep]
\item[(a)] $\psi$ is purely EM-sourced (role~2 only): the EM-sourced $\dpsi$
  may not produce gravitational-type refractive effects. The timing signal
  may not exist.
\item[(b)] $\psi$ is purely gravitational (role~1 only): EM sources do not
  generate $\dpsi$. The P1 experiment cannot create a psi-bubble.
\item[(c)] $\psi$ has both roles via a unified coupling not yet formulated:
  the current DFD theory is incomplete, and the P1 predictions are
  provisional until the unified coupling is derived.
\end{enumerate}

In all three cases, the P1 experimental predictions are either invalid
or rest on an incomplete theoretical foundation.
\end{theorem}

%=============================================================================
\section{Q4: Corrected Smoking-Gun Specification}
\label{sec:corrected_smoking_gun}
%=============================================================================

\subsection{Status of the W3i4 smoking-gun}
\label{ssec:smoking_gun_status}

The W3i4 smoking-gun specification (Table~4 of W3i4-P1) is
\textbf{physically unrealisable}:
\begin{itemize}[nosep]
\item It requires $|\dpsi| = 0.30$, which demands
  $u_{\rm EM} = 5.84\ee{16}$~J/m$^3$ ($>10^{12}\times$ any lab source).
\item The predicted signal (4.0~ns) and its discriminants (ON/OFF, rotation,
  sign reversal) are mathematically correct but refer to an unreachable
  operating point.
\end{itemize}

\subsection{Can any P1-type experiment serve as a smoking gun?}
\label{ssec:any_smoking_gun}

For a smoking gun, we need $\SNR \geq 5$ in a reasonable measurement time.
The signal scales as:
\begin{equation}
\delta t_{\rm NR} = \frac{2\,\lbone\,u_{\rm EM}\,R_b^2}{c^3} \times L/R_b
= \frac{2\,\lbone\,u_{\rm EM}\,R_b\,L}{c^3}.
\label{eq:signal_general}
\end{equation}

Maximising over all laboratory parameters with $L = R_b$:
\begin{equation}
\delta t_{\rm NR}^{(\rm max)} = \frac{2\,\lbone\,E_{\rm total}}{c^3}
\label{eq:signal_max}
\end{equation}
where $E_{\rm total} = u_{\rm EM}\,V$ is the total EM energy in the source.

For SNR $= 5$ with $\sigma_{\rm clock} = 1$~ps:
\begin{equation}
\lbone \geq \frac{5\,\sigma_{\rm clock}\,c^3}{2\,E_{\rm total}}
= \frac{5 \times 10^{-12} \times 2.7\ee{25}}{2\,E_{\rm total}}
= \frac{6.75\ee{13}}{E_{\rm total}}.
\label{eq:lambda_min_honest}
\end{equation}

\begin{center}
\textbf{Table 3: Minimum Detectable $\lbone$ vs.\ Source Energy}
\vspace{0.5em}

\begin{tabular}{@{}lcc@{}}
\toprule
\textbf{Source} & $E_{\rm total}$ (J) & $\lbone_{\rm min}$ (SNR $= 5$) \\
\midrule
TESLA SRF (1 kW) & 45.5 & $1.5\ee{12}$ \\
TESLA SRF (10 kW) & 909 & $7.4\ee{10}$ \\
TESLA SRF (100 kW) & 9090 & $7.4\ee{9}$ \\
LHC dipole circuit & $10^{10}$ & $6.75\ee{3}$ \\
Hypothetical $10^{15}$~J & $10^{15}$ & $6.75\ee{-2}$ \\
\midrule
\multicolumn{2}{@{}l}{\textbf{Current passive bound:}} & $1.56\ee{-9}$ \\
\bottomrule
\end{tabular}
\end{center}

\begin{theorem}[P1 Timing Experiment Cannot Reach Any Viable $\lbone$]
\label{thm:unreachable}
To detect a P1-type nonreciprocal timing signal at $\lbone = 1.56\ee{-9}$
(current passive bound), the required total EM source energy is:
\begin{equation}
E_{\rm min} = \frac{6.75\ee{13}}{1.56\ee{-9}} = 4.3\ee{22}~\text{J}
\approx 10^{16}~\text{kWh}.
\label{eq:E_min_passive}
\end{equation}

This is $\sim 10^{12}\times$ the LHC stored energy and $\sim 4\times$ the
annual global energy production ($\sim 6\ee{20}$~J). \textbf{No P1-type
timing experiment can reach the current passive bound on $\lbone$.}

Even at $\lbone = 10^{-2}$ (excluded): $E_{\rm min} = 6.75\ee{15}$~J
$= 1.88\ee{9}$~kWh $\approx 2$~GWh (the output of a nuclear power plant
for one hour, concentrated in a single EM field). While conceivable in
principle, this is $\sim 10^5\times$ the LHC and far beyond any planned
facility.
\end{theorem}

\subsection{Corrected smoking-gun: honest specification}
\label{ssec:corrected_spec}

Given the above, we provide an \emph{honest} corrected smoking-gun
specification:

\begin{tcolorbox}[colback=red!3!white, colframe=red!50!black,
  title=\textbf{Corrected Smoking-Gun Assessment for P1 (Field Drag / Gravitational Cloaking)}]

\textbf{Observable}: Nonreciprocal one-way timing asymmetry
$\mathcal{A}_{\rm NR} = t_{A\to B} - t_{B\to A}$ across an EM-generated
psi-field region.

\textbf{Signal strength at best achievable parameters}
(TESLA SRF, 10~kW, $\lbone = 1.56\ee{-9}$):
\begin{equation}
\delta t_{\rm NR} = 2.24\ee{-27}~\text{s} = 2.24~\text{ys (yoctoseconds)}.
\end{equation}

\textbf{Noise floor}: $\sigma_{\rm clock} = 1$~ps.

\textbf{Gap}: $\delta t_{\rm NR} / \sigma_{\rm clock} = 2.2\ee{-15}$
(15 orders of magnitude below detection).

\textbf{Verdict}: \textcolor{falsified}{\textbf{The P1 nonreciprocal timing
experiment is not a viable smoking gun for DFD.}} The signal-to-noise gap
of $10^{15}$ cannot be bridged by any combination of:
\begin{itemize}[nosep]
\item Source power increase (limited by SRF engineering to $\sim 10^4$~J stored).
\item Multi-patent enhancement (W3i4 Gamma configuration: $\sim 10^3\times$).
\item Measurement averaging (would require $>10^{22}$ repetitions).
\item Clock improvement (quantum limit already $\sim 1$~as, gains $\sim 10^3$).
\end{itemize}
Combined: $\sim 10^3 \times 10^3 \times 10^{11} = 10^{17}$ (the $10^{11}$
from $10^{22}$ repetitions at 1/s for $3\ee{14}$~years). Still $\sim 2$
orders short, and requiring $10^7$~years of averaging.

\textbf{Recommendation}: P1's smoking-gun role must be transferred to P11
(cavity mode detection) or P15 (direct psi-field generation/detection),
which detect resonant mode amplitudes where the signal scales differently.
However, the W3i4-P11 M$_\psi$ correction has degraded those estimates by
$10^6\times$, so those channels must also be re-evaluated at W3i5 level.
\end{tcolorbox}

%=============================================================================
\section{Impact of T5 on the P1 Source Equation}
\label{sec:T5_source}
%=============================================================================

\subsection{Which role does P1 require?}
\label{ssec:P1_role}

P1 requires \emph{both} roles of the psi-field simultaneously:
\begin{enumerate}[nosep]
\item \textbf{EM sourcing} (role~2): The SRF cavity generates $\dpsi$ via
  $\Box\dpsi = -(\lam-1)\,u_{\rm EM}/c^2$. Without this, no psi-bubble.
\item \textbf{Gravitational effect} (role~1): The generated $\dpsi$ modifies
  the spacetime metric as $g_{\mu\nu} \propto e^{\pm\dpsi}$, causing the
  refractive-index gradient that produces the timing signal. Without this,
  the psi-field is inert (no observable consequence for light propagation).
\end{enumerate}

\begin{finding}[P1 Is the Patent Most Vulnerable to T5]
\label{find:P1_T5_vulnerable}
Patent~1 is the only patent that requires \emph{both} the EM-sourcing and
the gravitational-metric roles of $\psi$ simultaneously in the same
experimental setup. Other patents either:
\begin{itemize}[nosep]
\item Use only the EM-sourcing role (P11, P15: detect the psi-mode amplitude
  via SQUID coupling, without needing $\psi$ to curve spacetime).
\item Use only the gravitational role (P5, P9, P13: apply DFD corrections to
  gravitational timing/navigation, without generating $\psi$ in the lab).
\end{itemize}

P1's drag-reduction and gravitational-cloaking claims are therefore the
\emph{most speculative} of all DFD patents, as they rest on the unresolved
T5 inconsistency.
\end{finding}

%=============================================================================
\section{Salvage Analysis: Is There Any P1-Adjacent Prediction That Survives?}
\label{sec:salvage}
%=============================================================================

\subsection{K-H instability suppression}
\label{ssec:KH}

The W3i1 finding that K-H suppression has a threshold five orders below
drag reduction ($\lbone_{\rm KH} \sim 10^{-7}$) remains a valid relative
statement. However, the same energy-density argument applies:
\begin{equation}
|\dpsi|_{\rm KH} \sim 10^{-8} \implies u_{\rm EM}^{(\rm KH)}
= \frac{10^{-8}\,c^2}{\lbone\,R_b^2}
= \frac{10^{-8} \times 9\ee{16}}{10^{-7} \times 25}
= 3.6\ee{14}~\text{J/m$^3$}.
\label{eq:KH_energy}
\end{equation}

This is still $\sim 10^{10}\times$ the SRF cavity capability.
K-H suppression is equally unachievable.

\subsection{Astrophysical observations}
\label{ssec:astro}

The P1-related predictions that \emph{do not} require laboratory generation
of $\dpsi$ include:
\begin{enumerate}[nosep]
\item NS timing anomaly: 2.15\% deviation from GR, $\delta t = 6.0$~ps (W3i4
  correction). Requires sub-ps pulsar timing (SKA may approach this in the
  2030s). This is a \emph{natural} psi-field from the NS gravitational
  potential, not a laboratory-generated one.
\item Gravitational lensing anomalies: the DFD refractive index $n = e^{\dpsi}$
  differs from the GR deflection angle by $\order{|\dpsi|^2}$ at the lens.
  For galaxy-cluster lenses ($|\dpsi| \sim 10^{-5}$), the deviation is
  $\order{10^{-10}}$ --- far below current lensing measurement precision
  ($\sim 1\%$).
\end{enumerate}

\begin{finding}[Only the NS Timing Anomaly Survives as a Testable P1-Adjacent Prediction]
\label{find:NS_survives}
Of all P1-related predictions, only the neutron-star timing anomaly (6.0~ps
DFD-vs-GR deviation) is both:
(a)~generated by a natural (not laboratory) psi-source, and
(b)~potentially within reach of next-generation instruments (SKA).

However, this prediction also depends on role~1 of the psi-field
(gravitational potential), which is under tension from T5. If the DFD
psi-field is not the Newtonian potential, the NS timing anomaly prediction
is void.
\end{finding}

%=============================================================================
\section{W3i5 P1 Summary and Honest Assessment}
\label{sec:summary}
%=============================================================================

\begin{tcolorbox}[colback=red!3!white, colframe=red!60!black,
  title=\textbf{W3i5 P1: Complete Honest Assessment}]
\textbf{Status of Patent 1 (Field Drag / Gravitational Cloaking) after W3i5:}

\begin{enumerate}[label=\textbf{\arabic*.}, itemsep=6pt]

\item \textbf{The 4.0~ns signal does not survive.}
The W3i4 signal of 4.0~ns assumed $|\dpsi| = 0.30$, which requires an EM
energy density of $5.8\ee{16}$~J/m$^3$ --- exceeding the best SRF cavity
by $10^{12}\times$. The actual signal at achievable energy densities is
$\sim 10^{-27}$~s (yoctoseconds), which is $10^{15}\times$ below the noise
floor.

\item \textbf{M$_\psi$ correction: indirect impact only.}
The $10^6\times$ M$_\psi$ over-optimism from W3i4-P11 does not directly
enter the P1 signal formula (which depends on static $\dpsi$, not mode
dynamics). However, it triggered the re-examination that revealed the far
more serious energy-density gap.

\item \textbf{No experimental configuration can achieve detection.}
Even Config~Gamma (all seven non-redundant patents, \$50M) achieves
$\SNR = 6.8\ee{-7}$ at the actual $|\dpsi|$. The LHC stored energy
(10~GJ) would give $\SNR \sim 10^{-5}$ at $\lbone = 10^{-2}$.

\item \textbf{T1 falsification does not directly kill P1}, but removes
the astrophysical motivation: if the cosmological sector is falsified,
the MOND fit that constrains $\mu_0$ is unreliable, and $\lbone$ has
no astrophysical anchor.

\item \textbf{T5 inconsistency is devastating for P1.}
P1 uniquely requires both EM-sourcing and gravitational-metric roles of
$\psi$ simultaneously. The T5 finding that these roles are inconsistent
means P1's predictions rest on an unresolved theoretical foundation.

\item \textbf{Only the NS timing anomaly (6.0~ps) survives} as a testable
P1-adjacent prediction, but it depends on T5 resolution and requires
SKA-level pulsar timing.

\item \textbf{Drag reduction and gravitational cloaking are unfeasible.}
The drag threshold ($\lbone > 1.85\ee{-2}$) requires nuclear-scale EM
energy densities. Gravitational cloaking (hiding from gravitational
detectors) requires even higher fields. These are not engineering
challenges; they are fundamental physics limitations.
\end{enumerate}

\textbf{Recommendation}: Patent~1 should be reclassified from
``near-term experimental target'' to ``deep-future theoretical concept
contingent on T5 resolution and a breakthrough in EM energy density
exceeding current capabilities by $>10^{10}\times$.''
\end{tcolorbox}

%=============================================================================
\section{W3i5 P1 Corrections to W3i4}
\label{sec:corrections}
%=============================================================================

\begin{center}
\textbf{Table 4: W3i5 Corrections and New Findings}
\vspace{0.5em}

\begin{tabular}{@{}p{5cm}ccp{4cm}@{}}
\toprule
\textbf{Item} & \textbf{W3i4 Value} & \textbf{W3i5 Value} & \textbf{Status} \\
\midrule
$|\dpsi|_{\rm bubble}$ at SRF operating point
  & 0.30 & $6.1\ee{-14}$ & \textcolor{falsified}{\textbf{CRITICAL CORRECTION}} \\
$\delta t_{\rm NR}$ (signal) & 4.0~ns & 0.81~zs & \textcolor{falsified}{\textbf{CRITICAL CORRECTION}} \\
$\SNR_{\rm baseline}$ & 4000 & $7.9\ee{-10}$ & \textcolor{falsified}{\textbf{CRITICAL CORRECTION}} \\
Smoking-gun viability & Yes & \textbf{No} & \textcolor{falsified}{\textbf{INVALIDATED}} \\
$\SNR_{\rm Alpha}$ & $1.3\ee{5}$ & $2.6\ee{-8}$ & Rescaled \\
$\SNR_{\rm Beta}$ & $1.1\ee{5}$ & $2.2\ee{-8}$ & Rescaled \\
$\SNR_{\rm Gamma}$ & $3.4\ee{6}$ & $6.8\ee{-7}$ & Rescaled \\
P1 patent classification & Near-term target & Deep-future concept & Downgraded \\
\midrule
\multicolumn{4}{@{}l}{\textbf{Unaffected W3i4 results:}} \\
\midrule
DFD PPN: $\beta = \gamma = 1$ & \checkmark & \checkmark & Confirmed \\
NS timing anomaly: 6.0~ps & \checkmark & \checkmark & Confirmed \\
Noise budget (96.1\% clock) & \checkmark & \checkmark & Confirmed \\
15-patent relative gains $G_k$ & \checkmark & \checkmark & Confirmed \\
Non-reciprocal drag geometry & \checkmark & \checkmark & Confirmed \\
\bottomrule
\end{tabular}
\end{center}

%=============================================================================
\section{Open Questions for W3i6}
\label{sec:open}
%=============================================================================

\begin{enumerate}[label=\textbf{OQ-\arabic*.}, itemsep=4pt]
\item \textbf{T5 resolution}: Can the DFD field equation be extended to
  include both gravitational and EM sourcing consistently? This is the
  existential question for P1.
\item \textbf{Alternative detection channels}: If P1 timing is unfeasible,
  which P11/P15 cavity-detection channels survive the M$_\psi$ correction?
  (Deferred to W3i5-P11 agent.)
\item \textbf{NS timing at SKA}: What is the realistic timeline and
  sensitivity for detecting 6.0~ps deviations in millisecond pulsar
  timing with SKA Phase~2?
\item \textbf{Energy-density scaling}: Is there a regime (e.g.,
  ultrashort-pulse lasers with coherent superposition) where $u_{\rm EM}$
  can exceed $10^{10}$~J/m$^3$ in a volume $>10^{-6}$~m$^3$ for durations
  $> 1$~ns?
\item \textbf{Cross-check with P2}: Does the P2 resonator analysis
  (toroidal re-entrant cavity) suffer from the same energy-density gap?
  The P2 ``$Q_\psi = 10^4$'' enhancement assumes a psi-mode resonance;
  does this require achievable $u_{\rm EM}$?
\end{enumerate}

%=============================================================================
% END W3i5 P1 UPDATE
%=============================================================================
